% TOP_PROBLEM: {"id": 3232, "title": "Partial List Coloring", "category_id": 3, "category": {"id": 3, "name": "graph_theory", "display_name": "Graph Theory", "description": "Problems involving graphs, networks, and their properties.", "slug": "graph-theory", "order_index": 3, "created_at": "2026-07-31T15:26:25.671Z"}, "status": "open", "problem_number": "OPG-788", "set_id": 12}
% FIRST_POSTED: 2026-10-01
% ENTRY_KIND: statement_only
% UnsolvedMath ID 3232; source catalogue code OPG-788
% !TeX program = lualatex
% Definition and sources only; this file is not an LLM solution attempt.
\documentclass[11pt,a4paper]{article}
\usepackage[margin=25mm]{geometry}
\usepackage{fontspec}
\setmainfont{lmroman10-regular.otf}[BoldFont=lmroman10-bold.otf,
 ItalicFont=lmroman10-italic.otf,BoldItalicFont=lmroman10-bolditalic.otf]
\usepackage{amsmath,amssymb,mathtools}
\usepackage{unicode-math}
\setmathfont{latinmodern-math.otf}
\AtBeginDocument{\let\setminus\smallsetminus}
\usepackage{xurl,hyperref}
\hypersetup{colorlinks=true,urlcolor=blue}
\setcounter{secnumdepth}{0}
\setlength{\parindent}{0pt}
\setlength{\parskip}{5pt}
\setlength{\emergencystretch}{3em}
\begin{document}
\section*{Partial List Coloring}\label{3232}
{\small UnsolvedMath ID 3232 · OPG-788 · Graph Theory · OpenGarden}
\subsection{Definitions and mathematical statement}
Let $G=(V,E)$ be a finite nonempty simple undirected graph. A list assignment
$L$ associates a finite set $L(v)$ of available colours with each vertex $v$.
An $L$-colouring of an induced subgraph $G[U]$, where $U\subseteq V$, is a
map $c:U\to\bigcup_{v\in V}L(v)$ such that $c(v)\in L(v)$ and
$c(u)\ne c(v)$ whenever $uv\in E$ has both ends in $U$. Put
\[
 \lambda_L(G)=\max\{|U|:G[U]\text{ has an }L\text{-colouring}\},
 \qquad
 \lambda_t(G)=\min_{|L(v)|=t\ (v\in V)}\lambda_L(G).
\]
The minimum ranges over arbitrary colour sets and arbitrary overlaps between
lists. Although there are infinitely many names for colours, the possible
values are integers between $0$ and $|V|$, so this minimum is well defined.

The list chromatic number $\chi_\ell(G)$ is the least positive integer $s$
for which every assignment of $s$ colours per vertex colours all of $G$.
The conjecture asks whether, for every such graph and every pair of integers
$1\le r\le s\le\chi_\ell(G)$,
\[
                 \frac{\lambda_r(G)}r\ge\frac{\lambda_s(G)}s.
\]
The quantity compares the guaranteed number of colourable vertices per
available colour. The optimizing subset may depend on the entire list
assignment; no particular subset must work for every assignment. At the
endpoint $s=\chi_\ell(G)$ the assertion reads
$\lambda_r(G)\ge r|V|/\chi_\ell(G)$, since all vertices can then be coloured.
\subsection{Short English statement}
Does the worst-case number of vertices that can be list-coloured, divided by the list size, decrease as the lists become larger?
\subsection{Sources}
\begin{itemize}
\item[S1] ULAM.AI and the UnsolvedMath Contributors, supplied problems.json, record 3232 (OPG-788), OpenGarden collection. Catalogue identity and source transcription; CC BY 4.0.
\url{https://huggingface.co/datasets/ulamai/UnsolvedMath}.
\item[S2] Open Problem Garden, Partial List Coloring. Original problem and discussion; the mathematical formulation below is expanded from the supplied source record.
\url{http://www.openproblemgarden.org/op/partial_list_coloring_0}.
\end{itemize}
\end{document}
